\subsection{Independent Ising constructions used for validation}
\label{sec:ising-independent}
The production reduced states use the exact fermion covariance method
of Sec.~\ref{sec:ising-production}. The full-vector constructions below
provide independent checks of that input: conformal blocks for $I,\eps$,
and conformal-block-seeded Lanczos projection for $\sigma$.

The chiral Ising fields are $1,\chi,\sigma$ with weights $0,1/2,1/16$;
$\chi$ is a Majorana fermion. Their fusion rules are
\begin{equation}
 \sigma\times\sigma=1+\chi,\qquad
 \chi\times\chi=1,\qquad \chi\times\sigma=\sigma.
\end{equation}

\subsubsection{Vacuum block: fusion data become spin amplitudes}
For the vacuum, the auxiliary correlator is
$\langle0|\sigma(z_1)\cdots\sigma(z_{2N})|0\rangle_{\bm p}$.
The subscript $\bm p$ chooses its intermediate fusion channels.
\begin{enumerate}
\item Use $2N$ spin fields on the circle,
$z_r=\exp[i\theta_r]$ with $\theta_r=\pi r/N$, $r=1,\ldots,2N$.
Pair $(z_{2j-1},z_{2j})$. Each pair fuses to identity ($m_j=0$) or
fermion ($m_j=1$), defining a physical spin with $Z_j=(-1)^{m_j}$.
\item Label the successive intermediate channels by bits
$\bm p=(p_1,\ldots,p_{N-1})$, setting $p_0=p_N=0$.
The physical bits are $m_j=p_{j-1}\mathbin{\mathrm{xor}}p_j$.
Thus $\sum_jm_j$ is even. The construction spans the even sector of the
physical spin parity $P=\prod_j Z_j$.
\item Enumerate binary strings $\bm q=(q_1,\ldots,q_{N-1})$ and define
$s_1=1$, $s_j=\prod_{k<j}(1-2q_k)$.
Split the field coordinates into two groups
$l_{\bm q}(j)=2j-(1+s_j)/2$ and
$l'_{\bm q}(j)=2j-(1-s_j)/2$. Each group contains one field from each pair.
Calculate the positive weight
\begin{equation}
 A_{\bm q}=\prod_{j>k}
 \left[\sin\frac{\theta_{l_{\bm q}(j)}-\theta_{l_{\bm q}(k)}}2
 \sin\frac{\theta_{l'_{\bm q}(j)}-\theta_{l'_{\bm q}(k)}}2\right]^{1/2}.
\end{equation}
\item Define $B_{\bm p}=\sum_{\bm q}(-1)^{\bm p\cdot\bm q}A_{\bm q}$,
where the dot product is a sum of products of bits. The vacuum amplitudes
are $\psi_I(\bm m(\bm p))\propto\sqrt{B_{\bm p}}$.
The binary sign sum is a Walsh--Hadamard transform, evaluated with
$O(N2^{N-1})$ butterfly operations instead of evaluating every
$(\bm p,\bm q)$ pair separately.
The weights themselves cost $O(N^2 2^{N-1})$ operations.
\item Normalize by
$\mathcal N_I=(\sum_{\bm p}B_{\bm p})^{1/2}$, so that
$\psi_I(\bm m(\bm p))=\sqrt{B_{\bm p}}/\mathcal N_I$.
Set all odd-parity amplitudes to zero. The output is a normalized
$2^N$-component vector.
\end{enumerate}
The homogeneous choice sets the pair-displacement parameter $\delta=0$ in
\texttt{fusion\_blocks}. Ref.~\cite{Montes2017} proves that this block
vacuum is the finite-chain critical ground state. We additionally check
its energy, symmetry sector, and residual numerically.

\subsubsection{Energy primary: insert two Majorana endpoint fields}
\begin{enumerate}
\item Keep the same $2N$ spin fields, coordinates, and physical basis.
Replace the auxiliary vacuum states by Majorana states:
\begin{equation}
 \Psi_\eps(\bm p)\propto
 \langle\chi|\sigma(z_1)\cdots\sigma(z_{2N})|\chi\rangle_{\bm p}.
\end{equation}
Since the chiral Majorana has weight $1/2$, the insertion at infinity
means $\chi(\infty)=\lim_{W\to\infty}W\chi(W)$.
The two endpoints preserve even physical parity.
\item Reuse $A_{\bm q}$ and $B_{\bm p}$ from the vacuum construction.
The endpoint correlator introduces the real factor
$C_{\bm q}=\cos[(\pi/2N)\sum_j s_j]$. Compute a second sign transform
$U_{\bm p}=\sum_{\bm q}(-1)^{\bm p\cdot\bm q}A_{\bm q}C_{\bm q}$.
The explicit result is
\begin{equation}
 \psi_\eps(\bm m(\bm p))\propto
 \frac{\displaystyle\sum_{\bm q}(-1)^{\bm p\cdot\bm q}A_{\bm q}
 \cos\left(\frac\pi{2N}\sum_{j=1}^Ns_j\right)}{\sqrt{B_{\bm p}}}.
\label{eq:epsblock}
\end{equation}
\item Normalize by
$\mathcal N_\eps=(\sum_{\bm p}|U_{\bm p}|^2/B_{\bm p})^{1/2}$
and embed $\psi_\eps(\bm m(\bm p))=U_{\bm p}/(\sqrt{B_{\bm p}}\mathcal N_\eps)$
in the same spin basis.
\item Check its orthogonality to $I$, its even parity, zero momentum,
energy, and eigenstate residual. The identification with the nonchiral
primary $\eps$ uses these spectral and symmetry properties; the
auxiliary chiral endpoint fields alone do not establish that identification.
\end{enumerate}
Equation~\eqref{eq:epsblock} is the published expression in
Ref.~\cite{Montes2017}, Eqs.~(73)--(75), specialized to the uniform
circle. Its equality with the desired excited eigenstate is numerically
verified here through $N=24$; we do not replace that numerical evidence
by an unprovided general proof for arbitrary insertions.

\subsubsection{Spin primary: conformal-block-seeded Lanczos}
\label{sec:sigmamethod}
An even-parity block cannot represent the odd-parity $\sigma$ state.
Ref.~\cite{Montes2017}, Sec.~X, reports that the naive endpoint-spin and
single-fermion candidates do not directly yield the desired homogeneous
odd-sector ground state. To check the production covariance state
independently at small sizes, we construct a full-vector representative
of $\sigma$ by a Hamiltonian projection initialized from the block vacuum.
This seed and projection are our combination of the published block
vacuum and a standard eigensolver, not a formula claimed in that paper.
Here ``block-seeded'' means seeded by a conformal block; the solver is
single-vector Lanczos, not a multi-vector block-Lanczos algorithm.
The steps below specify the numerical construction, including the part
performed by the Lanczos solver~\cite{ARPACK,ScipyEigsh}.

\begin{enumerate}
\item \textbf{Prepare a seed with the required symmetry.}
Starting from normalized $\ket{I_{\rm CB}}$, set
\begin{equation}
 \ket{v_\sigma}=\sum_{j=0}^{N-1}X_j\ket{I_{\rm CB}},\qquad
 \ket{q_1}=\frac{\ket{v_\sigma}}{\|v_\sigma\|_2}.
 \label{eq:sigmaseed}
\end{equation}
Since $PX_j=-X_jP$, the seed is odd. The sum is invariant under
translation, so the seed has zero momentum. It generally contains
higher-energy odd states as well as $\sigma$; applying the summed spin
operator is not itself the final projection.
\item \textbf{Build the odd-sector Hamiltonian action.}
Retain only bit strings $\bm m$ with $\sum_jm_j$ odd, leaving
$2^{N-1}$ basis vectors. Let $H_-$ denote $H$ restricted to this basis.
For a coefficient vector $f$, its action is
\begin{equation}
 (H_-f)_{\bm m}=
 -\left[\sum_j(-1)^{m_j}\right]f_{\bm m}
 -\sum_j f_{\bm m\oplus\bm e_j\oplus\bm e_{j+1}},
 \label{eq:oddaction}
\end{equation}
where $\bm e_j$ has one bit at site $j$, $\oplus$ is bitwise exclusive-or,
and site indices are periodic. The first term is diagonal; the second
flips neighboring bits. No full $2^{N-1}$ by $2^{N-1}$ matrix is stored.
Only parity is explicitly imposed in this basis; translation symmetry
is inherited from the seed and preserved by the Hamiltonian.
\item \textbf{Generate a Krylov space.}
After $r$ basis vectors, the search space is
\begin{equation}
 \mathcal K_r=\operatorname{span}
 \{q_1,H_-q_1,\ldots,H_-^{r-1}q_1\}.
\end{equation}
Lanczos constructs orthonormal vectors $q_j$. In exact arithmetic,
with $q_0=0$ and $b_0=0$, the recurrence is
\begin{align}
 a_j&=\langle q_j|H_-|q_j\rangle,\notag\\
 \ket{w_j}&=H_-\ket{q_j}-a_j\ket{q_j}-b_{j-1}\ket{q_{j-1}},\notag\\
 b_j&=\|w_j\|_2,\qquad \ket{q_{j+1}}=\ket{w_j}/b_j.
 \label{eq:lanczos}
\end{align}
If $b_j=0$, the generated subspace is already invariant. In floating-point
arithmetic the solver also controls loss of orthogonality. The symbols
$a_j,b_j$ here are Lanczos coefficients, unrelated to the boson endpoint
labels or the leading fitting exponent $\alpha$.
\item \textbf{Select the lowest-energy combination.}
Form $(T_r)_{ij}=\langle q_i|H_-|q_j\rangle$. This small matrix is
tridiagonal, with diagonal $a_j$ and neighboring entries $b_j$.
Let $\theta_{\min}$ and $\bm y$ be its smallest eigenvalue and normalized
eigenvector. The Rayleigh--Ritz approximation is
\begin{equation}
 \ket{\sigma_N^{(r)}}=\sum_{j=1}^r y_j\ket{q_j}.
\end{equation}
It minimizes the energy within the current search space. Increasing the
space, or restarting while retaining the useful spectral information,
refines the approximation.
\item \textbf{Converge and measure the actual residual.}
The code calls \texttt{eigsh} with \texttt{k=1}, \texttt{which="SA"},
the seed as \texttt{v0}, and \texttt{tol=2e-12}. It uses implicitly
restarted Lanczos~\cite{ScipyEigsh}. The requested solver tolerance is
not an absolute bound on $\eta$: after convergence we independently
evaluate Eq.~\eqref{eq:eigenresidual} using the original Hamiltonian.
\item \textbf{Identify and validate the primary.}
Normalize the vector, restore zero entries on even bit strings, and
check $P=-1$, zero momentum, orthogonality to $I,\eps$, and the energy
$E_\sigma=-2\cot[\pi/(2N)]$. The vanishing large-$N$ gap has scaling
dimension $1/8$, identifying the $(1/16,1/16)$ primary.
Our explicit construction checks use
$e_P=\|(P+1)\sigma_N\|_2$ and
$e_{\mathsf T}=\|(\mathsf T-1)\sigma_N\|_2$ for the two symmetry errors.
\end{enumerate}

The spectral projection underlying these steps can also be written as
\begin{equation}
 \ket{\sigma_N}=\lim_{\tau\to\infty}
 \frac{e^{-\tau(H-E_*)}\ket{q_1}}
 {\|e^{-\tau(H-E_*)}\ket{q_1}\|_2}.
 \label{eq:projection}
\end{equation}
Here $\tau$ is imaginary time and $E_*$ is an arbitrary reference energy.
The normalized seed has an expansion into odd-sector eigenstates;
provided its $\sigma$ coefficient is nonzero, higher energies are
suppressed relative to the lowest one. Lanczos implements the search
without explicitly exponentiating $H$.
The measured squared seed overlap
$F_{\rm seed}=|\langle\sigma_N|q_1\rangle|^2$ is approximately
$0.986,0.983,0.981,0.979,0.977$ at $N=6,8,10,12,16$.
At $N=16$, the recorded projected-state residual is
$2.70\times10^{-12}$.
The complete size-by-size residual and independent-state checks appear
in Appendix~\ref{sec:ising-spin-validation}; the largest recorded residual is
$1.99\times10^{-11}$ at $N=12$. Thus the solver tolerance must not be
quoted as the achieved absolute residual.

This method obtains the required lattice primary using both a conformal
block and the Hamiltonian. A pure endpoint conformal-block formula for
the periodic $\sigma$ state remains unestablished here.

